\begin{textsample}{POS Dim 6 – global\_north\_2024\_05 – Score 10.00 – global\_north\_2024\_05/108528869.txt}  \label{ex:f6_pos_144}
US-built \textbf{pier} will be removed from Gaza \textbf{coast} and repaired after damage from rough \textbf{seas} A U.S.
Army landing craft is seen beached in Ashdod on Sunday, May 26, 2024, after being swept by wind and current from the temporary humanitarian \textbf{pier} in the Gaza Strip. ( AP Photo/Tsafrir Abayov ) WASHINGTON ( AP ) -- The U.S.-built temporary \textbf{pier} that has been taking humanitarian \textbf{aid} to \textbf{starving} Palestinians for less than two weeks will be removed from the \textbf{coast} of Gaza to be repaired after getting damaged in rough \textbf{seas} and weather, the Pentagon said Tuesday.
Over the next two days, the \textbf{pier} will be pulled from the beach and sent to the southern Israeli city of Ashdod, where U.S.
Central Command will repair it, Pentagon spokeswoman Sabrina Singh told reporters.
She said the fixes will take"at least over a week"and then the \textbf{pier} will need to be anchored back into the beach in Gaza."From when it was operational, it was working, and we just had sort of an unfortunate confluence of weather storms @ @ @ @ @ @ @ @ @ @ said."Hopefully just a little over a week, we should be back up and running."The \textbf{pier}, used to carry in humanitarian \textbf{aid} arriving by \textbf{sea}, is one of the few ways that free food and other supplies are getting to Palestinians who the U.N. says are on the brink of \textbf{famine} amid the nearly 8-month-old war between Israel and Hamas in Gaza.
The two main crossings in southern Gaza, Rafah from Egypt and Kerem Shalom from Israel, are either not operating or are largely inaccessible for the U.N. because of fighting nearby as Israel pushes into Rafah.
The \textbf{pier} and two crossings from Israel in northern Gaza are where most of the incoming humanitarian \textbf{aid} has entered in the past three weeks.
The setback is the latest for the \$320 million \textbf{pier}, which only began operations in the past two weeks and has already had three U.S. service members injured and had four \textbf{vessels} beached due to heavy \textbf{seas}.
Two of the service members received minor injuries but the third @ @ @ @ @ @ @ @ @ @ Deliveries also were halted for two days last week after crowds rushed \textbf{aid} trucks coming from the \textbf{pier} and one Palestinian man was shot dead.
The \textbf{pier} was fully functional as late as Saturday when heavy \textbf{seas} unmoored four of the Army \textbf{boats} that were being used to ferry pallets of \textbf{aid} from commercial \textbf{vessels} to the \textbf{pier}.
The system is anchored into the beach in Gaza and provided a long causeway for trucks to drive that \textbf{aid} onto the shore.
Two of the \textbf{vessels} got stuck on the \textbf{coast} of Israel.
One has already been recovered and the other will be in the next 24 hours with the help of the Israeli military, Singh said.
The other two \textbf{boats} were stranded on the beach in Gaza and were expected to be recovered in the next two days, she said.
U.S. officials have repeatedly emphasized that the \textbf{pier} can not provide the amount of \textbf{aid} that \textbf{starving} Gazans need and said that more checkpoints for humanitarian trucks need to be opened.
At maximum capacity, the \textbf{pier} would @ @ @ @ @ @ @ @ @ @, and U.S. officials have stressed the need for open land crossings for the remaining 1.8 million.
The U.S. also has planned to continue to provide \textbf{airdrops} of food, which likewise can not meet all the needs.
A deepening Israeli offensive in the southern city of Rafah has made it impossible for \textbf{aid} \textbf{shipments} to get through the crossing there, which is a key source for fuel and food coming into Gaza.
Israel says it is bringing \textbf{aid} in through another border crossing, Kerem Shalom, but humanitarian organizations say Israeli military operations make it difficult for them to retrieve the \textbf{aid} there for distribution.

% matched lemmas: aid, airdrop, boat, coast, famine, pier, sea, shipment, starve, vessel
\end{textsample}
